%% appendix-search.tex — Search Protocol and Included Studies
%% Included via %% appendix-search.tex — Search Protocol and Included Studies
%% Included via %% appendix-search.tex — Search Protocol and Included Studies
%% Included via %% appendix-search.tex — Search Protocol and Included Studies
%% Included via \input{appendix-search} at the end of sn-article.tex.

\section{Search protocol}
\label{app:search}

\subsection*{A1.\ Sources, strings and record flow}

All searches were run on \cutoffdate{} for records published from January~2022.
The date filter was applied in the query for OpenAlex (\texttt{from\_publication\_date})
and Scopus (\texttt{PUBYEAR > 2021}), and after retrieval for arXiv, OpenReview and
IEEE~Xplore. Duplicates were identified first by arXiv identifier or DOI and then by
normalised title (Unicode NFKD, ASCII folding, case and punctuation removed).
Table~\ref{tab:flow} gives the record flow; the search scripts are
available at the project repository; the exact query strings as executed are
listed below.

\begin{table}[h]
\caption{Record flow of the structured search.}
\label{tab:flow}
\small
\begin{tabular}{lrrr}
\toprule
Source & Retrieved & Outside window & In window \\
\midrule
OpenAlex    & \NretOpenAlex{}   & \NoutOpenAlex{}   & \NinOpenAlex{} \\
Scopus      & \NretScopus{}     & \NoutScopus{}     & \NinScopus{} \\
arXiv       & \NretarXiv{}      & \NoutarXiv{}      & \NinarXiv{} \\
OpenReview  & \NretOpenReview{} & \NoutOpenReview{} & \NinOpenReview{} \\
IEEE Xplore & \NretIEEE{}       & \NoutIEEE{}       & \NinIEEE{} \\
\midrule
Total       & \NrecTotal{}      & \Noutwindow{}     & \NinTotal{} \\
\multicolumn{3}{l}{Duplicates removed} & \Ndup{} \\
\multicolumn{3}{l}{Unique records screened (title/abstract)} & \Nscreened{} \\
\multicolumn{3}{l}{Full texts assessed} & \Nfulltext{} \\
\multicolumn{3}{l}{Evidence studies included via the structured search} & \NincludedS{} \\
\bottomrule
\end{tabular}
\end{table}

\smallskip
\paragraph{Recursive-training terms (TERMS).}
``model collapse'' | ``knowledge collapse'' | ``self-consuming'' |
``model autophagy'' | ``curse of recursion'' | autophagous | ``recursive training'' |
``iterative retraining'' | ``self-trained'' | ``data feedback loop(s)'' |
``closed-loop learning'' | ``self-correcting loop'' | ``iterated learning'' |
``self-rewarding'' | ``iterative DPO'' | ``iterative preference''.

\paragraph{Generative-model terms (GEN).}
``generative model'' | ``language model'' | LLM | ``diffusion model'' |
``foundation model'' | ``synthetic data'' | ``synthetic text'' | ``generated data''.

\paragraph{Per-source syntax.}
\emph{OpenAlex}: TERMS AND GEN in title and abstract;
\texttt{from\_publication\_date:2022-01-01}; cursor-paginated.
\emph{Scopus}: \texttt{TITLE-ABS-KEY(TERMS) AND TITLE-ABS-KEY(GEN)};
wildcard \texttt{self-train*}; \texttt{PUBYEAR > 2021}; institutional access CAPES/UFPA.
\emph{arXiv}: each TERMS phrase in \texttt{ti:} or \texttt{abs:}; GEN required in \texttt{abs:};
categories cs.LG, cs.CL, cs.CR, cs.AI; 54 pre-2022 records excluded after retrieval.
\emph{OpenReview}: \NqOpenReview{} phrase queries; the query list omits
``model autophagy'', ``curse of recursion'', ``autophagous'' and ``self-trained''
because the OpenReview interface does not support them reliably;
year derived from a truncated Unix timestamp, decoded post-retrieval.
\emph{IEEE~Xplore}: \NqIEEE{} paired queries (a TERMS phrase paired with a GEN phrase
where the interface required it); date filter post-retrieval.

\smallskip
\noindent\textit{Semantic Scholar:} not searched owing to API rate limits at the cutoff date.

% ─────────────────────────────────────────────────────────────────────────────

\subsection*{A2.\ List of Included Sources}

The manuscript cites \Ncited~sources in total: \Nevidence~evidence studies,
\Nposition~position paper(s), \Ncontext~context sources, and \Nlegal~legal
and normative instruments (\SuppStudies).
Their category (E~evidence, P~position, C~context, L~legal/normative), entry route
(S~structured search, A~anchor, CH~citation chasing,
R~added during revision to address a reviewer comment or to verify a claim),
anchor status, and preprint/workshop flag (\dag) are provided in
\SuppStudies.
Evidence studies ($N{=}\Nevidence{}$) are those whose formal results, experiments or
empirical observations are directly synthesised in the body of the review;
position, context and legal/normative sources are cited for arguments, background or regulatory reference.
Category reflects how each source is used in this review.
 at the end of sn-article.tex.

\section{Search protocol}
\label{app:search}

\subsection*{A1.\ Sources, strings and record flow}

All searches were run on \cutoffdate{} for records published from January~2022.
The date filter was applied in the query for OpenAlex (\texttt{from\_publication\_date})
and Scopus (\texttt{PUBYEAR > 2021}), and after retrieval for arXiv, OpenReview and
IEEE~Xplore. Duplicates were identified first by arXiv identifier or DOI and then by
normalised title (Unicode NFKD, ASCII folding, case and punctuation removed).
Table~\ref{tab:flow} gives the record flow; the search scripts are
available at the project repository; the exact query strings as executed are
listed below.

\begin{table}[h]
\caption{Record flow of the structured search.}
\label{tab:flow}
\small
\begin{tabular}{lrrr}
\toprule
Source & Retrieved & Outside window & In window \\
\midrule
OpenAlex    & \NretOpenAlex{}   & \NoutOpenAlex{}   & \NinOpenAlex{} \\
Scopus      & \NretScopus{}     & \NoutScopus{}     & \NinScopus{} \\
arXiv       & \NretarXiv{}      & \NoutarXiv{}      & \NinarXiv{} \\
OpenReview  & \NretOpenReview{} & \NoutOpenReview{} & \NinOpenReview{} \\
IEEE Xplore & \NretIEEE{}       & \NoutIEEE{}       & \NinIEEE{} \\
\midrule
Total       & \NrecTotal{}      & \Noutwindow{}     & \NinTotal{} \\
\multicolumn{3}{l}{Duplicates removed} & \Ndup{} \\
\multicolumn{3}{l}{Unique records screened (title/abstract)} & \Nscreened{} \\
\multicolumn{3}{l}{Full texts assessed} & \Nfulltext{} \\
\multicolumn{3}{l}{Evidence studies included via the structured search} & \NincludedS{} \\
\bottomrule
\end{tabular}
\end{table}

\smallskip
\paragraph{Recursive-training terms (TERMS).}
``model collapse'' | ``knowledge collapse'' | ``self-consuming'' |
``model autophagy'' | ``curse of recursion'' | autophagous | ``recursive training'' |
``iterative retraining'' | ``self-trained'' | ``data feedback loop(s)'' |
``closed-loop learning'' | ``self-correcting loop'' | ``iterated learning'' |
``self-rewarding'' | ``iterative DPO'' | ``iterative preference''.

\paragraph{Generative-model terms (GEN).}
``generative model'' | ``language model'' | LLM | ``diffusion model'' |
``foundation model'' | ``synthetic data'' | ``synthetic text'' | ``generated data''.

\paragraph{Per-source syntax.}
\emph{OpenAlex}: TERMS AND GEN in title and abstract;
\texttt{from\_publication\_date:2022-01-01}; cursor-paginated.
\emph{Scopus}: \texttt{TITLE-ABS-KEY(TERMS) AND TITLE-ABS-KEY(GEN)};
wildcard \texttt{self-train*}; \texttt{PUBYEAR > 2021}; institutional access CAPES/UFPA.
\emph{arXiv}: each TERMS phrase in \texttt{ti:} or \texttt{abs:}; GEN required in \texttt{abs:};
categories cs.LG, cs.CL, cs.CR, cs.AI; 54 pre-2022 records excluded after retrieval.
\emph{OpenReview}: \NqOpenReview{} phrase queries; the query list omits
``model autophagy'', ``curse of recursion'', ``autophagous'' and ``self-trained''
because the OpenReview interface does not support them reliably;
year derived from a truncated Unix timestamp, decoded post-retrieval.
\emph{IEEE~Xplore}: \NqIEEE{} paired queries (a TERMS phrase paired with a GEN phrase
where the interface required it); date filter post-retrieval.

\smallskip
\noindent\textit{Semantic Scholar:} not searched owing to API rate limits at the cutoff date.

% ─────────────────────────────────────────────────────────────────────────────

\subsection*{A2.\ List of Included Sources}

The manuscript cites \Ncited~sources in total: \Nevidence~evidence studies,
\Nposition~position paper(s), \Ncontext~context sources, and \Nlegal~legal
and normative instruments (\SuppStudies).
Their category (E~evidence, P~position, C~context, L~legal/normative), entry route
(S~structured search, A~anchor, CH~citation chasing,
R~added during revision to address a reviewer comment or to verify a claim),
anchor status, and preprint/workshop flag (\dag) are provided in
\SuppStudies.
Evidence studies ($N{=}\Nevidence{}$) are those whose formal results, experiments or
empirical observations are directly synthesised in the body of the review;
position, context and legal/normative sources are cited for arguments, background or regulatory reference.
Category reflects how each source is used in this review.
 at the end of sn-article.tex.

\section{Search protocol}
\label{app:search}

\subsection*{A1.\ Sources, strings and record flow}

All searches were run on \cutoffdate{} for records published from January~2022.
The date filter was applied in the query for OpenAlex (\texttt{from\_publication\_date})
and Scopus (\texttt{PUBYEAR > 2021}), and after retrieval for arXiv, OpenReview and
IEEE~Xplore. Duplicates were identified first by arXiv identifier or DOI and then by
normalised title (Unicode NFKD, ASCII folding, case and punctuation removed).
Table~\ref{tab:flow} gives the record flow; the search scripts are
available at the project repository; the exact query strings as executed are
listed below.

\begin{table}[h]
\caption{Record flow of the structured search.}
\label{tab:flow}
\small
\begin{tabular}{lrrr}
\toprule
Source & Retrieved & Outside window & In window \\
\midrule
OpenAlex    & \NretOpenAlex{}   & \NoutOpenAlex{}   & \NinOpenAlex{} \\
Scopus      & \NretScopus{}     & \NoutScopus{}     & \NinScopus{} \\
arXiv       & \NretarXiv{}      & \NoutarXiv{}      & \NinarXiv{} \\
OpenReview  & \NretOpenReview{} & \NoutOpenReview{} & \NinOpenReview{} \\
IEEE Xplore & \NretIEEE{}       & \NoutIEEE{}       & \NinIEEE{} \\
\midrule
Total       & \NrecTotal{}      & \Noutwindow{}     & \NinTotal{} \\
\multicolumn{3}{l}{Duplicates removed} & \Ndup{} \\
\multicolumn{3}{l}{Unique records screened (title/abstract)} & \Nscreened{} \\
\multicolumn{3}{l}{Full texts assessed} & \Nfulltext{} \\
\multicolumn{3}{l}{Evidence studies included via the structured search} & \NincludedS{} \\
\bottomrule
\end{tabular}
\end{table}

\smallskip
\paragraph{Recursive-training terms (TERMS).}
``model collapse'' | ``knowledge collapse'' | ``self-consuming'' |
``model autophagy'' | ``curse of recursion'' | autophagous | ``recursive training'' |
``iterative retraining'' | ``self-trained'' | ``data feedback loop(s)'' |
``closed-loop learning'' | ``self-correcting loop'' | ``iterated learning'' |
``self-rewarding'' | ``iterative DPO'' | ``iterative preference''.

\paragraph{Generative-model terms (GEN).}
``generative model'' | ``language model'' | LLM | ``diffusion model'' |
``foundation model'' | ``synthetic data'' | ``synthetic text'' | ``generated data''.

\paragraph{Per-source syntax.}
\emph{OpenAlex}: TERMS AND GEN in title and abstract;
\texttt{from\_publication\_date:2022-01-01}; cursor-paginated.
\emph{Scopus}: \texttt{TITLE-ABS-KEY(TERMS) AND TITLE-ABS-KEY(GEN)};
wildcard \texttt{self-train*}; \texttt{PUBYEAR > 2021}; institutional access CAPES/UFPA.
\emph{arXiv}: each TERMS phrase in \texttt{ti:} or \texttt{abs:}; GEN required in \texttt{abs:};
categories cs.LG, cs.CL, cs.CR, cs.AI; 54 pre-2022 records excluded after retrieval.
\emph{OpenReview}: \NqOpenReview{} phrase queries; the query list omits
``model autophagy'', ``curse of recursion'', ``autophagous'' and ``self-trained''
because the OpenReview interface does not support them reliably;
year derived from a truncated Unix timestamp, decoded post-retrieval.
\emph{IEEE~Xplore}: \NqIEEE{} paired queries (a TERMS phrase paired with a GEN phrase
where the interface required it); date filter post-retrieval.

\smallskip
\noindent\textit{Semantic Scholar:} not searched owing to API rate limits at the cutoff date.

% ─────────────────────────────────────────────────────────────────────────────

\subsection*{A2.\ List of Included Sources}

The manuscript cites \Ncited~sources in total: \Nevidence~evidence studies,
\Nposition~position paper(s), \Ncontext~context sources, and \Nlegal~legal
and normative instruments (\SuppStudies).
Their category (E~evidence, P~position, C~context, L~legal/normative), entry route
(S~structured search, A~anchor, CH~citation chasing,
R~added during revision to address a reviewer comment or to verify a claim),
anchor status, and preprint/workshop flag (\dag) are provided in
\SuppStudies.
Evidence studies ($N{=}\Nevidence{}$) are those whose formal results, experiments or
empirical observations are directly synthesised in the body of the review;
position, context and legal/normative sources are cited for arguments, background or regulatory reference.
Category reflects how each source is used in this review.
 at the end of sn-article.tex.

\section{Search protocol}
\label{app:search}

\subsection*{A1.\ Sources, strings and record flow}

All searches were run on \cutoffdate{} for records published from January~2022.
The date filter was applied in the query for OpenAlex (\texttt{from\_publication\_date})
and Scopus (\texttt{PUBYEAR > 2021}), and after retrieval for arXiv, OpenReview and
IEEE~Xplore. Duplicates were identified first by arXiv identifier or DOI and then by
normalised title (Unicode NFKD, ASCII folding, case and punctuation removed).
Table~\ref{tab:flow} gives the record flow; the search scripts are
available at the project repository; the exact query strings as executed are
listed below.

\begin{table}[h]
\caption{Record flow of the structured search.}
\label{tab:flow}
\small
\begin{tabular}{lrrr}
\toprule
Source & Retrieved & Outside window & In window \\
\midrule
OpenAlex    & \NretOpenAlex{}   & \NoutOpenAlex{}   & \NinOpenAlex{} \\
Scopus      & \NretScopus{}     & \NoutScopus{}     & \NinScopus{} \\
arXiv       & \NretarXiv{}      & \NoutarXiv{}      & \NinarXiv{} \\
OpenReview  & \NretOpenReview{} & \NoutOpenReview{} & \NinOpenReview{} \\
IEEE Xplore & \NretIEEE{}       & \NoutIEEE{}       & \NinIEEE{} \\
\midrule
Total       & \NrecTotal{}      & \Noutwindow{}     & \NinTotal{} \\
\multicolumn{3}{l}{Duplicates removed} & \Ndup{} \\
\multicolumn{3}{l}{Unique records screened (title/abstract)} & \Nscreened{} \\
\multicolumn{3}{l}{Full texts assessed} & \Nfulltext{} \\
\multicolumn{3}{l}{Evidence studies included via the structured search} & \NincludedS{} \\
\bottomrule
\end{tabular}
\end{table}

\smallskip
\paragraph{Recursive-training terms (TERMS).}
``model collapse'' | ``knowledge collapse'' | ``self-consuming'' |
``model autophagy'' | ``curse of recursion'' | autophagous | ``recursive training'' |
``iterative retraining'' | ``self-trained'' | ``data feedback loop(s)'' |
``closed-loop learning'' | ``self-correcting loop'' | ``iterated learning'' |
``self-rewarding'' | ``iterative DPO'' | ``iterative preference''.

\paragraph{Generative-model terms (GEN).}
``generative model'' | ``language model'' | LLM | ``diffusion model'' |
``foundation model'' | ``synthetic data'' | ``synthetic text'' | ``generated data''.

\paragraph{Per-source syntax.}
\emph{OpenAlex}: TERMS AND GEN in title and abstract;
\texttt{from\_publication\_date:2022-01-01}; cursor-paginated.
\emph{Scopus}: \texttt{TITLE-ABS-KEY(TERMS) AND TITLE-ABS-KEY(GEN)};
wildcard \texttt{self-train*}; \texttt{PUBYEAR > 2021}; institutional access CAPES/UFPA.
\emph{arXiv}: each TERMS phrase in \texttt{ti:} or \texttt{abs:}; GEN required in \texttt{abs:};
categories cs.LG, cs.CL, cs.CR, cs.AI; 54 pre-2022 records excluded after retrieval.
\emph{OpenReview}: \NqOpenReview{} phrase queries; the query list omits
``model autophagy'', ``curse of recursion'', ``autophagous'' and ``self-trained''
because the OpenReview interface does not support them reliably;
year derived from a truncated Unix timestamp, decoded post-retrieval.
\emph{IEEE~Xplore}: \NqIEEE{} paired queries (a TERMS phrase paired with a GEN phrase
where the interface required it); date filter post-retrieval.

\smallskip
\noindent\textit{Semantic Scholar:} not searched owing to API rate limits at the cutoff date.

% ─────────────────────────────────────────────────────────────────────────────

\subsection*{A2.\ List of Included Sources}

The manuscript cites \Ncited~sources in total: \Nevidence~evidence studies,
\Nposition~position paper(s), \Ncontext~context sources, and \Nlegal~legal
and normative instruments (\SuppStudies).
Their category (E~evidence, P~position, C~context, L~legal/normative), entry route
(S~structured search, A~anchor, CH~citation chasing,
R~added during revision to address a reviewer comment or to verify a claim),
anchor status, and preprint/workshop flag (\dag) are provided in
\SuppStudies.
Evidence studies ($N{=}\Nevidence{}$) are those whose formal results, experiments or
empirical observations are directly synthesised in the body of the review;
position, context and legal/normative sources are cited for arguments, background or regulatory reference.
Category reflects how each source is used in this review.
